\documentclass[14pt,aspectratio=1610]{beamer}
%\documentclass[14pt,aspectratio=1610,handout]{beamer}
\usepackage{csu-theme}
\usepackage{binary-trees}
\usepackage{arrays}

\title{\Huge{} Binary Heaps \& Heapsort}
\subtitle{CSCI 315: \textit{Data Structure Analysis}\\Chapter 10 of \href{https://opendatastructures.org/ods-cpp.pdf}{Open Data Structures}}
%\author{Sean T. Hayes, Ph.D.}
\institute{Charleston Southern University}

%% Comment out date to use default (today)
% \date{March 18, 2025}
\subject{Software Programming}
%\keywords{}

\setlength{\parskip}{0.5em}
\tikzset{
array index/.style = {
  font=\small\ttfamily\bfseries
},
}

\begin{document}

\begin{frame}
  \titlepage{}
\end{frame}

\section*{Why Heapsort?}

\begin{frame}{Why Binary Heaps?}

The \emph{binary heap} is a fundamental structure used\\ in many algorithms.

\begin{itemize}

  \item  A heap supports very efficient operations:

  \begin{itemize}
    \item Insert: \(O(\log_2{n})\)

    \item Extract max/min: \(O(\log_2{n})\)

    \item Peek/Poll max/min: \(O(1)\)
  \end{itemize}

  This structure underpins algorithms like:

  \begin{itemize}
    \item Dijkstra's Shortest Path Algorithm

    \item Prim's Minimum-Spanning Tree Algorithm

    \item Priority queues used in schedulers and simulations.\\
      Will improve our \emph{Priority Queue} design.
  \end{itemize}

  \item Understanding heapsort means you understand how heaps work internally, not just how to call a library.
\end{itemize}
\end{frame}

\begin{frame}{Why Heapsort? --- Guaranteed Performance}
  \begin{itemize}
\item Heapsort guarantees \(O(n \log n)\) time complexity.
\item This bound holds in \textbf{best, average, and worst cases}.
\item Comparison with other sorting algorithms:
\begin{itemize}
\item Quicksort: average \(O(n \log n)\) but worst case \(O(n^2)\)
\item Merge Sort: \(O(n \log n)\) but requires extra memory for arrays
\item Heapsort: \(O(n \log n)\) with guaranteed worst-case performance
\end{itemize}
\item This makes heapsort useful when \textbf{predictable performance} is required.
\end{itemize}
\end{frame}

\begin{frame}{Why Heapsort? --- Memory Efficiency}
\begin{itemize}
\item Heapsort is an \textbf{in-place sorting algorithm}.
\item It requires only \(O(1)\) additional memory.
\item Comparison:
\begin{itemize}
\item Heapsort: \(O(1)\) auxiliary space
\item Merge Sort: \(O(n)\) auxiliary space
\end{itemize}
\item This property is useful in:
\begin{itemize}
\item Memory-constrained systems
\item Embedded systems
\item Large datasets where extra memory is expensive
\end{itemize}
\end{itemize}
\end{frame}

\begin{frame}{Why Heapsort? --- Educational Value}
\begin{itemize}
\item Reinforces key computer science concepts:
\begin{itemize}
\item Binary trees
\item Priority queues
\item Time complexity analysis
\item Array-based tree representations
\end{itemize}
\item Common interview topics include:
\begin{itemize}
\item Building a heap
\item Heapify operations
\item Finding the $k$ largest or smallest elements
\end{itemize}
\item Understanding heapsort demonstrates strong algorithmic foundations.
\end{itemize}
\end{frame}

% \begin{frame}{Why Heapsort?}
% \begin{itemize}[<+->]
% \item
%   Heapsort is a well-known, traditional sorting\\
%   algorithm you are expected to know.
% \item
%   Heapsort is \emph{always} $O(n \log n)$.
% \item
%   Learn about binary \textit{search} trees to:
%   \begin{itemize}
%   \item
%     Quicksort is usually $O(n \log n)$ but $O(n^2)$ in the worst case.
%   \item
%     Quicksort is sometimes faster, but Heapsort is better in time-critical applications.
%   \item
%     Merge Sort is also $O(n \log n)$, but requires additional memory to sort an array.
%   \end{itemize}
% \item
%   Can improve our Priority Queue design (in certain situations).
% \end{itemize}
% \end{frame}

\section{Heaps}

\begin{frame}{What is a ``heap''?}
\begin{itemize}[<+->]
\item
  There are two definitions of \emph{heap}.
  \begin{enumerate}
  \item
    \emph{Memory}: A large area of memory from which the \\
    programmer can allocate blocks as needed, and \\
    deallocate them (or allow them to be garbage collected) when no longer needed.
  \item
    \emph{Data Structure}: A balanced, left-justified binary tree in which no node has a value greater than its parent (binary heap).
  \end{enumerate}
\item
  These two definitions have little in common.
\item
  Heapsort uses the second definition.
\end{itemize}
\end{frame}

\begin{frame}{Two Properties of a Binary Heap}
The heap data structure is a binary tree that maintains\\
two properties, which we will discuss.\\[1em]
\begin{enumerate}
\item
  \emph{Shape Property}: \textit{Left justified} and, therefore, \textit{balanced}
\item
  \emph{Heap Property}: The relationship between elements
\end{enumerate}
\end{frame}

\begin{frame}{Balanced Binary Trees}
\begin{itemize}
\item<2->
  Recall:
  \begin{itemize}
  \item
    The \emph{depth of a node} is its distance from the root.
  \item<3->
    The \emph{depth of a tree} is the depth of the deepest node.
  \end{itemize}
\item<4->
  A binary tree of depth \(n\) is \emph{balanced} if all the nodes \\
  at depths 0 through $n-2$ have two children.
\end{itemize}
\begin{columns}[T]
\column{0.4\textwidth}%
\onslide<5->{%
  \begin{forest}
    [, for tree={fit=tight, l=7mm, inner sep=0.2em, s sep-=1mm}, name=root
      [, name=lev1
        [, name=lev2
          [, name=levN]
          []
        ]
        [
          []
          []
        ]
      ]
      [
        [
          []
          []
        ]
        [
          []
          []
        ]
      ]
    ]
    \coordinate (labelLine) at ([yshift=0.25em]root);
    \path (root |- labelLine) node[anchor=south,align=center] {(a) Balanced};
    \onslide<6->{ % Show levels
      \coordinate (levelLine) at ([left,xshift=-1mm]levN);
      \path (levelLine |- root) node[anchor=east, align=right] {\footnotesize $0$};
      \path (levelLine |- lev1) node[anchor=east, align=right] {\footnotesize $n-2$};
      \path (levelLine |- lev2) node[anchor=east, align=right] {\footnotesize $n-1$};
      \path (levelLine |- levN) node[anchor=east, align=right] {\footnotesize \(n\)};
    }
  \end{forest}%
}

\column{0.3\textwidth}%
\only<7->{%
  \begin{forest}
    [, for tree={fit=tight, l=7mm, inner sep=0.2em, s sep-=1mm}, name=root
      [
        [
          []
          []
        ]
        [
          [, phantom]
          [, phantom]
        ]
      ]
      [
        [
          [, phantom]
          []
        ]
        [
          []
          [, phantom]
        ]
      ]
    ]
    \coordinate (labelLine) at ([yshift=0.25em]root);
    \path (root |- labelLine) node[anchor=south,align=center] {(b) Balanced};
  \end{forest}%
}

\column{0.3\textwidth}%
\only<8->{%
  \begin{forest}
    [, for tree={fit=tight, l=7mm, inner sep=0.2em, s sep-=1mm}, name=root
      [
        [
          []
          []
        ]
        [
          []
          []
        ]
      ]
      [
        [
          []
          []
        ]
        [, phantom]
      ]
    ]
    \coordinate (labelLine) at ([yshift=0.25em]root);
    \path (root |- labelLine) node[anchor=south,align=center] {(c) Unbalanced};
  \end{forest}%
}
\end{columns}
\end{frame}



\begin{frame}{Shape Property: Left-Justified Binary Trees}
A binary tree has the \emph{shape properly} if:
\begin{itemize}
\item<2->
  Balanced
\item<3->
  All leaves at depth \(n\) are to the left, without any missing.\\
  (All leaves at depth \(n - 1\) are to the right of all the nodes at depth \(n\).)
\end{itemize}
\begin{columns}[T]
\centering
\onslide<4->{\column{0.35\textwidth}\centering
  \begin{forest}
    [, for tree={fit=tight, l=7mm, inner sep=0.25em}, name=root
      [
        [
          []
          []
        ]
        [
          []
          []
        ]
      ]
      [
        [
          []
          [, phantom]
        ]
        [
          [, phantom]
          [, phantom]
        ]
      ]
    ]
    \coordinate (labelLine) at ([yshift=0.25em]root);
    \path (root |- labelLine) node[anchor=south,align=center] {(a) Left-Justified};
  \end{forest}%
}

\onslide<5->{\column{0.35\textwidth}\centering
  \begin{forest}
    [, for tree={fit=tight, l=7mm, inner sep=0.25em}, name=root
      [
        [
          []
          []
        ]
        [
          [, phantom]
          []
        ]
      ]
      [
        [
          []
          [, phantom]
        ]
        [
          []
          []
        ]
      ]
    ]
    \coordinate (labelLine) at ([yshift=0.25em]root);
    \path (root |- labelLine) node[anchor=south,align=center] {(b) Not Left-Justified};
  \end{forest}%
}
\end{columns}
\end{frame}

\begin{frame}{The Heap Property}
\begin{itemize}[<+->]
\item
  A node has the \emph{heap property} if the value in the\\
  node is at least as large as its children's values.
\item
  All leaves automatically have the heap property.
\item
  A binary tree is a \emph{heap} if all nodes in it have the heap property.
\end{itemize}

\begin{columns}[T]
\centering
\onslide<4->{\column{0.18\textwidth}\centering
  \begin{forest}
    for tree={fit=tight, minimum size=1.8em, l sep=1em, inner sep=0}
    [12,name=root, draw=CSUcyan,fill=white,
      [8]
      [3]
    ]
    \node[anchor=south,align=center,above=0.2em of root] {(a) heap};
  \end{forest}%
}

\onslide<5->{\column{0.18\textwidth}\centering
  \begin{forest}
    for tree={fit=tight, minimum size=1.8em, l sep=1em, inner sep=0}
    [12,name=root, draw=CSUcyan,fill=white,
      [8]
      [12]
    ]
    \node[anchor=south,align=center,above=0.2em of root] {(b) heap};
  \end{forest}%
}

\onslide<6->{\column{0.18\textwidth}\centering
  \begin{forest}
    for tree={fit=tight, minimum size=1.8em, l sep=1em, inner sep=0}
    [8,name=root, draw=CSUcyan,fill=white,
      [12]
      [3]
    ]
    \node[anchor=south,align=center,above=0.2em of root] {(c) not heap};
  \end{forest}%
}
\end{columns}
\end{frame}

\begin{frame}{Example Heap}
\centering
  \begin{forest}
    for tree={minimum size=1.8em, l sep=1em, inner sep=0}
    [14,name=root
      [12
        [5
          [1]
          [, phantom]
        ]
        [12]
      ]
      [4
        [2]
        [3]
      ]
    ]
  \end{forest}
\end{frame}

\begin{frame}{Shift Up}
If a node is without the heap property, exchange \\
its value with the larger child's value.

\begin{columns}[T]
\centering
\onslide<2->{\column{0.25\textwidth}\centering
  \begin{forest}
    for tree={fit=tight, minimum size=1.8em, l sep=1em, inner sep=0}
    [12,name=root, draw=CSUcyan,fill=white,
      [8]
      [14]{
        \onslide<3->{\draw[<->,dotted, very thick] () to[bend right] node[midway,right] {sift} (root);}
      }
    ]
    \node[anchor=south,align=center,above=0.2em of root] {(a) not heap};
  \end{forest}%
}

\onslide<4->{\column{0.25\textwidth}\centering
  \begin{forest}
    for tree={fit=tight, minimum size=1.8em, l sep=1em, inner sep=0}
    [14,name=root, draw=CSUcyan,fill=white,
      [8]
      [12]
    ]
    \node[anchor=south,align=center,above=0.2em of root] {(b) heap};
  \end{forest}%
}
\end{columns}

\begin{itemize}
\item<5->
  This is sometimes called \emph{sifting up} or bubble up.
\item<6->
  Notice that the child may have \emph{lost} the heap property during shifting.
\end{itemize}
\end{frame}

\begin{frame}{Plan of Attack}
We will discover:\\[1em]
\begin{enumerate}
\item<2->
  How to turn a binary tree into a heap.
\item<3->
  After particular changes, how to turn a binary tree \emph{back} into a heap.
\item<4->
  How to use these ideas to sort an \emph{array}.\\
  (This is the cool part.)
\end{enumerate}
\end{frame}

\section{Insert}

\begin{frame}{Inserting into a Heap}
\begin{itemize}[<+->]
\item
  A single-node tree is automatically a heap.
\item
  We construct a heap by adding nodes one at a time:
  \begin{itemize}
  \item
    Add the node just to the right of the rightmost node in the deepest level.
  \item
    If the deepest level is full, start a new level.
  \end{itemize}
\end{itemize}
\end{frame}

\begin{frame}{Order of Insertion}
  \centerline{%
    \begin{forest}
      [, name=root, for tree={fit=tight, l=8mm, inner sep=0.25em},
        for tree={visible on=<2->}
        [, for tree={visible on=<3->}
          [, for tree={visible on=<5->}
            [, for tree={visible on=<9->}]
            [, for tree={visible on=<10->}]
          ]
          [, for tree={visible on=<6->}
            [, for tree={visible on=<11->}]
            [, phantom]
          ]
        ]
        [, for tree={visible on=<4->}
          [, for tree={visible on=<7->}
            [, phantom]
            [, phantom]
          ]
          [, for tree={visible on=<8->}
            [, phantom]
            [, phantom]
          ]
        ]
      ]
  \end{forest}%
  }
\end{frame}

\begin{frame}{Inserting into a Heap}
\begin{itemize}[<+->]
\item
  Problem: Each time we add a node, we may\\ destroy the heap property of its parent node.
\item
  Solution: \emph{Sift up} (a.k.a., bubble up).
\item
  But each time we sift up, the value of the topmost node in the sift may increase, and this may destroy the heap property of its parent node.
\item
  We repeat the sifting up process, moving up in the tree, until either:
  \begin{itemize}
  \item
    We reach a node whose value doesn't need to be swapped (because the parent is larger than both children), or
  \item
    We reach the root.
  \end{itemize}
\end{itemize}
\end{frame}

\begin{frame}{Inserting into a Heap}
Insert 8, 10, 5, 11, 6, 14, and 12.
\begin{center}
  \begin{forest}
    for tree={fit=band, minimum size=1.8em, l sep=1em, inner sep=0}
    [\alt<-4>{8}{\alt<-11>{10}{\alt<-18>{11}{14}}}, name=root, for tree={visible on=<2->}
      [\alt<-4>{10}{\alt<-9>{8}{\alt<-11>{11}{10}}}, name=1_1, for tree={visible on=<3->}
        [\alt<-9>{11}{8}, for tree={visible on=<8->}]{
          \onslide<9-10>{\draw[<->,dotted, very thick] () to[bend left] node[midway,left] {sift} (1_1);}}
        [6, for tree={visible on=<14->}]
      ]{
        \onslide<4-5,11-12>{\draw[<->,dotted, very thick] () to[bend left] node[midway,left] {sift} (root);}}
      [\alt<-16>{5}{\alt<-18>{14}{\alt<-22>{11}{12}}}, name=1_2, for tree={visible on=<7->}
        [\alt<-16>{14}{5}, for tree={visible on=<15->}]{
          \onslide<16-17>{\draw[<->,dotted, very thick] () to[bend right] node[midway,right] {sift} (1_2);}}
        [\alt<-22>{12}{11}, visible on=<21->]{
          \onslide<22-23>{\draw[<->,dotted, very thick] () to[bend right] node[midway,right] {sift} (1_2);}}
      ]{
        \onslide<18-19>{\draw[<->,dotted, very thick] () to[bend right] node[midway,right] {sift} (root);}}
    ]\onslide<24>{}
  \end{forest}
\end{center}
\onslide<20->{}
\end{frame}

\begin{frame}{A Sample Heap}
\begin{columns}
\column{0.4\textwidth}
Here is a sample \emph{heapified} binary tree.
\begin{itemize}
\item<2->
  Heapified does \emph{not} mean sorted.
\item<3->
  Heapifying does \emph{not} change the shape of the binary tree.\\
  The binary tree was always balanced and left-justified.
\end{itemize}

\column{0.6\textwidth}
  \centerline{%
    \begin{forest}
      for tree={fit=band, l sep=1em, inner sep=0em, minimum size=1.8em}
      [25
        [22
          [19
            [18]
            [14]
          ]
          [22
            [21]
            [3]
          ]
        ]
        [17
          [14
            [9]
            [11]
          ]
          [15]
        ]
      ]
    \end{forest}%
  }
\end{columns}
\end{frame}


\section{Remove}

\begin{frame}[t]{Removing the Root}
\begin{columns}[T]
\column{0.4\textwidth}
\begin{itemize}
\item<2->
  The largest number is now in the root.
\item<3->
  Suppose we \emph{discard} the root.
\item<5->
  Problem: The tree becomes \emph{unbalanced} and \emph{un-left-justified}.
\item<6->
  Solution: Make the rightmost leaf the new root.
\end{itemize}
\column{0.6\textwidth}
\vspace{1.5cm}
  \centerline{%
    \begin{forest}
      for tree={fit=band, l sep=1em, inner sep=0em, minimum size=1.8em}
      [\alt<-7>{25}{11}, /tikz/invisible on=<5-7> %\alt<-3>{25}{\alt<-7>{}{11}}, name=root, justText on=<4-7>, minimum size=2.1em
        [22
          [19
            [18]
            [14]
          ]
          [22
            [21]
            [3]
          ]
        ]
        [17
          [14
            [9]
            [11, visible on=<-7>]{
              \onslide<7>{\draw[->,dotted, very thick] () to[bend right] node[midway,left] {replace} (root);}}
          ]
          [15]
        ]
      ]
      \onslide<8>{}% Make root show up again at the end
    \end{forest}%
  }
\end{columns}
\end{frame}

\begin{frame}[t]{Re-Heap by Sifting Down}
\begin{columns}[T]
  \column{0.4\textwidth}
  \begin{itemize}
  \item<2->
    Our tree is balanced and left-justified, but no longer a heap.
  \item<3->
    Only the root lacks the heap property.
  \item<4->
    We \emph{sift down} (or trickle down) the root.
  \item<8->
    Now, only one child may have lost the heap property.
  \end{itemize}
\column{0.6\textwidth}
\vspace{1.5cm}
  \centerline{%
    \begin{forest}
      for tree={fit=band, l sep=1em, inner sep=0em, minimum size=1.8em}
      [\alt<-5>{11}{22}, name=root
        [\alt<-5>{22}{11}, name=1_1
          [19
            [18]
            [14]
          ]
          [22
            [21]
            [3]
          ]
        ]{
          \onslide<5-6>{\draw[<->,dotted, very thick] () to[bend left] node[midway,left] {sift} (root);}}
        [17
          [14
            [9]
            [, phantom]
          ]
          [15]
        ]
      ]
    \end{forest}%
  }
\end{columns}
\end{frame}

\begin{frame}[t]{Re-Heap (Sift-Down)}
\begin{columns}[T]
  \column{0.4\textwidth}
  \begin{itemize}
  \item<2->
    Now, the left child of the root (currently 11) lacks the heap property.
  \item<3->
    We can sift down this node.
  \item<6->
    Now, only one of its children may have lost the heap property.
  \end{itemize}
\column{0.6\textwidth}
\vspace{1.5cm}
  \centerline{%
    \begin{forest}
      for tree={fit=band, l sep=1em, inner sep=0em, minimum size=1.8em}
      [22, name=root
        [\alt<-4>{11}{22}, name=1_1
          [19
            [18]
            [14]
          ]
          [\alt<-4>{22}{11}
            [21]
            [3]
          ]{
            \onslide<4-5>{\draw[<->,dotted, very thick] () to[bend right] node[midway,right] {sift} (1_1);}}
        ]
        [17
          [14
            [9]
            [, phantom]
          ]
          [15]
        ]
      ]
    \end{forest}%
  }
\end{columns}
\end{frame}


\begin{frame}[t]{Re-Heap (Sift-Down)}
\begin{columns}[T]
  \column{0.4\textwidth}
  \begin{itemize}
  \item<2->
    Now, the right child of the left child of the root (currently 11) lacks the heap property.
  \item<3->
    We can sift down this node.
  \item<6->
    It is a leaf, so there are children to lost the heap property.
  \end{itemize}
\column{0.6\textwidth}
\vspace{1.5cm}
  \centerline{%
    \begin{forest}
      for tree={fit=band, l sep=1em, inner sep=0em, minimum size=1.8em}
      [22, name=root
        [22, name=1_1
          [19
            [18]
            [14]
          ]
          [\alt<-4>{11}{21}, name=2_2
            [\alt<-4>{21}{11}]{
              \onslide<4-5>{\draw[<->,dotted, very thick] () to[bend left] node[midway,left] {sift} (2_2);}}
            [3]
          ]
        ]
        [17
          [14
            [9]
            [, phantom]
          ]
          [15]
        ]
      ]
    \end{forest}%
  }
\end{columns}
\end{frame}


\begin{frame}[t]{Re-Heap (Sift-Down)}
\begin{columns}[T]
\column{0.4\textwidth}
  \begin{itemize}[<+->]
  \item
    Every node again has the heap property.
  \item
    The largest value is at the root.
  \item
    Repeating until the tree is empty will produce produces a sequence ordered from largest to smallest.
  \end{itemize}
\column{0.6\textwidth}
\vspace{1.5cm}
  \centerline{%
    \begin{forest}
      for tree={fit=band, l sep=1em, inner sep=0em, minimum size=1.8em}
      [22, name=root
        [22, name=1_1
          [19
            [18]
            [14]
          ]
          [21, name=2_2
            [11]
            [3]
          ]
        ]
        [17
          [14
            [9]
            [, phantom]
          ]
          [15]
        ]
      ]
    \end{forest}%
  }
\end{columns}
\end{frame}

\section{Sorting}

% \defverbatim{\LstHeapPseudocode}{%
% \begin{lstlisting}[language={}]
% Heapify the array.
% While the array isn't empty:
% {
%   Remove and replace the root.
%   Reheap the new root node.
% }
% \end{lstlisting}}
\begin{frame}{Sorting}
\begin{columns}[T]
  \column{0.4\textwidth}
  What do heaps have to do with sorting an array?
  \begin{itemize}
  \item<2->
    Because the binary tree is balanced and left-justified, it can be represented as an array!
  \item<3->
    All our operations on binary trees can be represented as operations on arrays.
  \end{itemize}
\column{0.6\textwidth}
\only<4->{%
\begin{block}{Algorithm Pseudocode}\texttt{%
Heapify the array.\\
While the array isn't empty:\\
\{\\
\hspace{2em}Remove and replace the root.\\
\hspace{2em}Re-heap the new root node.\\
\}}
\end{block}%
}
\end{columns}
\end{frame}

\begin{frame}{Mapping into an Array}
\begin{columns}
  \column{0.43\textwidth}
    \begin{itemize}
    \item<3->
      The left child of index, $i$, is at $2 \times i + 1$
    \item<4->
      The right child of index, $i$, is at $2 \times i + 2$
    \item<5->
      Ex: Children of $3$ ($19$) are $7$ ($18$) and $8$ ($14$).
    \item<6->
      Given a child, what is it's parent's index?
    \end{itemize}
  \column{0.57\textwidth}
    \centerline{%
      \begin{forest}
        for tree={fit=band, l sep=1em, inner sep=0em, minimum size=1.8em}
        [25
          [22
            [19
              [18]
              [14]
            ]
            [22
              [21]
              [3]
            ]
          ]
          [17
            [14
              [9]
              [11]
            ]
            [15]
          ]
        ]
      \end{forest}%
    }
\end{columns}
\pause
\begin{tikzpicture}
  \ArrayLabel{}
  \ArrayList{25, 22, 17, 19, 22, 14, 15, 18, 14, 21, 3, 9, 11}
\end{tikzpicture}
\end{frame}

\begin{frame}[t]{Removing and Replacing the Root}
\vspace{2em}
\begin{tikzpicture}
  \ArrayLabel{}
  \alt<-5>{
    \ArrayList{25, 22, 17, 19, 22, 14, 15, 18, 14, 21, 3, 9, 11}
    \onslide<5> {%
      \draw[link, <->] ([xshift=-2mm] el25.south east) to[bend right=18] node[below] {swap} ([xshift=2mm] el11.south west);%
    }
  }{%
    \ArrayList{11, 22, 17, 19, 22, 14, 15, 18, 14, 21, 3, 9}
    \alt<-7> {
      \ArrayNode{25}
    }{
      \ArrayNode[MyGreen]{25}
    }

    \onslide<6> {%
      \draw[link, <->] ([xshift=-2mm] el11.south east) to[bend right=18] node[below] {swap} ([xshift=2mm] el25.south west);%
    }
  }
\end{tikzpicture}
\vspace{-1em}
\begin{itemize}
\item<2->
  The ``root'' is the first element in the array.
\item<3->
  The ``deepest rightmost node'' is the last element.
\item<4->
  Remove the root by swapping the first and last values.
\item<8->
  \ldots{}pretend that the last element in the array no \\longer exists --- the ``last index'' is 11 (9).
\end{itemize}

\end{frame}

\begin{frame}[t]{Reheap and Repeat}
  \vspace{2em}
\begin{minipage}[h][3cm][t]{\textwidth} % keeps array from shifting based on swap arrows
\begin{tikzpicture}
  \ArrayLabel{}
  \alt<-3>{%
    \ArrayNode{11} %
    \node[block] (root) [fit=(prev)]{};
    \ArrayNode{22} %
    \onslide<3>{\ArraySwap{el11}{el22}}%
  }{%
    \alt<-13>{
      \ArrayNode{22}
    }{
      \alt<-18>{
        \ArrayNode{9}
      }{
        \ArrayNode{22}
      }
    } %
    \node[block] (root) [fit=(prev)]{};
    \alt<-6>{
      \ArrayNode{11}
    }{
      \alt<-18>{
        \ArrayNode{22}
      }{
        \alt<-21>{
          \ArrayNode{9}
        }{
          \ArrayNode{21}
        }
      }
    } %
    \onslide<4>{\ArraySwap{el22}{el11}}%
    \only<18>{\ArraySwap{el9}{el22}}%
    \only<19>{\ArraySwap{el22}{el9}}%
  }

  \alt<-6>{
    \ArrayList{17, 19, 22, 14, 15, 18, 14, 21, 3}
    \onslide<6>{\ArraySwap{el11}{el22}}
  }{
    \alt<-9>{
      \ArrayList{17, 19, 11, 14, 15, 18, 14, 21, 3}
      \only<7>{\ArraySwap{el22}{el11}}
      \only<9>{\ArraySwap{el11}{el21}}
    }{
      \alt<-21>{
        \ArrayList{17, 19, 21, 14, 15, 18, 14, 11, 3}
      }{
        \alt<-24>{
          \ArrayList{17, 19, 9, 14, 15, 18, 14, 11, 3}
        }{
          \ArrayList{17, 19, 11, 14, 15, 18, 14, 9, 3}
        }
      }
      \only<10>{\ArraySwap{el21}{el11}}
      \only<21>{\ArraySwap{el9}{el21}}
      \only<22>{\ArraySwap{el21}{el9}}
      \only<24>{\ArraySwap{el9}{el11}}
      \only<25>{\ArraySwap{el11}{el9}}
    }
  }

  % Second-to-last element
  \alt<-13>{\ArrayNode{9}}{\alt<-15>{\ArrayNode{22}}{\ArrayNode[MyGreen]{22}}}
  \onslide<13-14>{\ArraySwap[18]{root}{prev}}

  \ArrayNode[MyGreen]{25} % Last element
  \onslide<26>{}
\end{tikzpicture}
\end{minipage}
\begin{itemize}
\item<2->
  \alt<-11>{Re-heap the root node (index 0, containing 11).}{
    Again, remove and replace the root node.
  }
\item<15->
  Remember, though, that the ``last'' array index is changed.
\item<17->
  Repeat until the last becomes the first, and the array is sorted!
\end{itemize}
\end{frame}

\section*{Implementation}

\begin{frame}{Analysis}\Large\centering
Let's start implementing a heap structure.
\end{frame}

\section{Analysis}

\begin{frame}{Analysis: Step 1, Heapify the Array}
The sorting algorithm starts by \emph{heapifying the array}.
\begin{itemize}[<+->]
\item
  We add each of \(n\) nodes.
  \begin{itemize}
  \item
    Each node must be sifted up, possibly as far as the root.\\
    Because the binary tree is perfectly balanced, sifting up a single node takes $O(\log n)$ time.
  \item
    Because we iterate \(n\) times, heapifying takes $n \times O(\log n)$ time, that is, $O(n \log n)$ time.
  \end{itemize}
\end{itemize}
\end{frame}

\begin{frame}{Analysis: Step 2, Remove}
\begin{itemize}
\item<1->
  Here's the rest of the algorithm:
  \begin{columns}\column{0.6\textwidth}
  \begin{block}{Algorithm Pseudocode (Part 2)}\texttt{%
    While the array isn't empty: \{\\
    \hspace{2em}Remove and replace the root.\\
    \hspace{2em}Reheap the new root node.\\
    \}}
  \end{block}
  \end{columns}

\item<2->
  We do the while loop \(n\) times (actually, $n-1$ times), because we remove one of the \(n\) nodes each time.
\item<3->
  Removing and replacing the root takes \(O(1)\) time.
\item<4->
  Therefore, the total time is \(n\) times however long\\ it takes the reheap method.
\end{itemize}
\end{frame}


\begin{frame}{Analysis: Step 3, Re-heap}
\begin{itemize}[<+->]
\item
  To reheap the root node, we follow {one path} \\
  from the root to a leaf (or until the heap \\
  property exists for the subtree).
\item
  The binary tree is perfectly balanced.
\item
  Therefore, this path is $O(\log n)$ long.
  \begin{itemize}
  \item
    And we only do \(O(1)\) operations at each node.
  \item
    Therefore, a reheap takes $O(\log n)$ time.
  \end{itemize}
\item
  Since we reheap inside a while loop that we do \(n\) times, the total time for the while loop is $n \times O(\log n)$, or $O(n \log n$).
\end{itemize}
\end{frame}

\begin{frame}{Analysis: Performance of All Steps}
\begin{columns}\column{0.6\textwidth}
  \begin{block}{Algorithm Pseudocode}\texttt{%
    Heapify the array.\\
    While the array isn't empty: \{\\
    \hspace{2em}Remove and replace the root.\\
    \hspace{2em}Reheap the new root node.\\
    \}}
  \end{block}
\end{columns}
\begin{itemize}
\item<+->
  We have seen that heapifying takes $O(n \log n)$ time.
\item<+->
  The while loop takes $O(n \log n)$ time.
\item<+->
  Therefore, the total time is $O(n \log n) + O(n \log n)$.
\item<+->
  Which is $O(n \log n)$ time (average- and worst-case).
\end{itemize}
\end{frame}

\begin{frame}{Summary of Benefits}
\begin{itemize}[<+->]
  \item
    \emph{Efficient Operations}:\\
    Insertion and deletion are $O(\log n)$.\\
    Retrieval of the maximum (or minimum) value is \(O(1)\).
  \item
    \emph{Space Efficiency}:\\
    Binary heaps can be implemented using arrays, removing the need for pointers.
  \item
    \emph{Ease of Implementation}:\\
    With arrays, the algorithms for insertion, deletion, and heapifying are intuitive.
  \end{itemize}
\end{frame}
\begin{frame}{Summary of Benefits}
  \begin{itemize}[<+->]
    \item
      \emph{Heapsort}: An always $O(n \log{n})$ in-place\\
      sorting algorithm\\
      (i.e., \(O(1)\) additional space and with predictable performance).
    \item
      \emph{Priority Queue}: Many algorithms require elements with higher (or lower) priority to be processed first.
\end{itemize}
\end{frame}

\begin{frame}{Visual Comparison of Sorting Algorithms}
  \begin{center}\Large
    \href{https://www.toptal.com/developers/sorting-algorithms}{Sorting Algorithms Animations}
  \end{center}
\end{frame}

\section{Introsort}

\begin{frame}{Introspective Sort (Introsort)}
  Combines \emph{Quicksort}, \emph{Insertion Sort}, and \emph{Heapsort},

  Approach:

  \begin{enumerate}
    \item 
      Start with \emph{Quicksort} for its fast average-case performance (\(\approx O(n \log_2 n)\)).

    \item 
      Track the recursion depth.

    \item 
      If the depth exceeds a threshold (typically \(2 \cdot{} \log_2 n\)), switch to \emph{Heapsort} to avoid quicksort's worst case (\(O(n^2)\)).

    \item 
      For small subarrays (size 16 to 32), use \emph{Insertion Sort} for efficiency.
  \end{enumerate}
\end{frame}


\section*{Lab 16}
\begin{frame}{Lab 16: Binary Heap and Heapsort}
\begin{center}\Large
Let's take a look at the lab based on today's lecture.
\end{center}
\end{frame}
\end{document}